\documentclass[11pt]{article}
\usepackage[margin=27mm]{geometry}
\usepackage{amsmath,amssymb,amsthm,booktabs,hyperref}
\newtheorem{lemma}{Lemma}
\newtheorem{theorem}[lemma]{Theorem}
\newcommand{\PSL}{\operatorname{PSL}}
\newcommand{\PGL}{\operatorname{PGL}}
\newcommand{\N}{\operatorname{N}}
\title{Counting configurations above a cyclic subgroup of order three}
\author{Human Compression Project}\date{}
\begin{document}\maketitle
\begin{abstract}
One regular action of a dihedral quotient counts all configurations
whose maximal parabolic is $A_4$ and whose intersection has order three.
Its reflections describe the second parabolics $S_3$, while its nonidentity
elements describe ordered pairs of $A_4$ parabolics. Passing to unordered
pairs identifies an element with its inverse. This recovers five rows
of the source table and corrects the conjugacy counts in its last two
rows. The full automorphism counts in those rows remain valid.
\end{abstract}

\section{The hypotheses and the common normalizer}
Let $p>3$ be prime such that $A_4$ is maximal in $G=\PSL(2,p)$.
The classical maximal-subgroup list gives exactly $p=5$ or
$p\equiv3,13,27,37\pmod {40}$. In particular $p\equiv\pm3\pmod8$.
Put $E=\PGL(2,p)=\operatorname{Aut}(G)$, and choose
$\epsilon\in\{1,-1\}$ with $3\mid p-\epsilon$. Write
\[
 n=\frac{p-\epsilon}{2}=3m,\qquad m=\frac{p-\epsilon}{6}.
\]
Fix $B\cong C_3$ in $G$. We retain the classical facts that all
such $B$ are conjugate, that the subgroups $A_4$ form one $G$-class,
and that $\N_E(A_4)\cong S_4$. The latter group acts on its normal
$A_4$ by the usual automorphisms. Also, the centralizer torus of $B$
has order $n$ in $G$ and $2n$ in $E$, and its normalizers are dihedral.
These are subgroup facts, not assumptions about the number of geometries.

Choose torus and reflection generators so that
\[
 \N_E(B)=\langle R,S\mid R^{2n}=S^2=1,\ SRS=R^{-1}\rangle,
 \quad \N_G(B)=\langle R^2,S\rangle,
 \quad B=\langle R^{2m}\rangle.
\]
After quotienting by $B$, write
\[
 \overline N=\N_E(B)/B=\langle x,s\mid x^{2m}=s^2=1,
 \ sxs=x^{-1}\rangle,
 \qquad K=\N_G(B)/B=\langle r=x^2,s\rangle.
\]
Thus $|K|=2m$ and $r$ has order $m$. This presentation also covers
$m=1,2$ without an exceptional convention for a small dihedral group.

\begin{lemma}\label{regular}
The set $\mathcal A$ of subgroups $A_4$ containing $B$ is a regular
$K$-set. Choose one of them, denoted $M$, as the identity point.
The stabilizer of $M$ in $\overline N$ is generated by an external
reflection $j$. Coordinates may be chosen with $j=xs$; its action
on the $K$-coordinates is the automorphism
\[
 \theta(r^k)=r^{-k},\qquad \theta(r^k s)=r^{1-k}s.
\]
\end{lemma}
\begin{proof}
All subgroups $C_3$ in $A_4$ are conjugate internally, and their
normalizer in $A_4$ is $C_3$. Any transporter between two members
of $\mathcal A$ can therefore be adjusted inside the target $A_4$
to normalize $B$. Hence $\N_G(B)$ is transitive on $\mathcal A$.
Since $M$ is maximal in the simple group $G$, it is self-normalizing.
The stabilizer of $M$ is consequently
$\N_G(B)\cap M=\N_M(B)=B$. Moreover $B$ normalizes every
member of $\mathcal A$, since it is contained in each. The action
therefore factors through $K$ and is regular.

The stabilizer in $\N_E(B)$ is $\N_{\N_E(M)}(B)$, the normalizer
of a three-cycle subgroup in $S_4$. It is $S_3$, with quotient $C_2$
by $B$. Its involution inverts $B$, so its image in $\overline N$
is a reflection. It is outside $K$, since the stabilizer in $K$ is
trivial. Every external reflection has the form $x^{2a+1}s$;
changing the initial reflection generator permits $j=xs$.
Conjugation by $j$ sends $r$ to $r^{-1}$ and $s$ to $rs$.
Since $j$ fixes the identity point $M$ in the regular set, these
conjugation formulas are exactly its action on the coordinates.
\end{proof}

\section{Different parabolic types}
The common overgroup calculation used in the source audit has the
following interface: if $H\le G$ properly contains $B$ and its action
on $H/B$ is two-transitive, then $H$ is $C_6$, $S_3$ or $A_4$.
For this particular $B$ the list can also be checked directly.
Writing $d=[H:B]$, transitivity of $B$ on the other $d-1$ cosets
gives $d-1\mid3$, so $d=2$ or $4$. At degree two $H$ has order
six and is cyclic or symmetric. At degree four the kernel, being
contained in $B$, has order one or three. A kernel of order three
would leave a group of order four acting two-transitively on four
points, impossible. Thus the action is faithful and its image is the
unique subgroup $A_4$ of order twelve in $S_4$. This proves the interface
without an additional finite search.

\begin{theorem}\label{different}
The configurations with parabolics $M\cong A_4$ and $H\cong C_6$
exist exactly when $m$ is even; then there is one conjugacy class
and one $E$-class. With $H\cong S_3$, there are $m$ conjugacy
classes and $\lceil m/2\rceil$ classes under $E$.
\end{theorem}
\begin{proof}
A cyclic overgroup $C_6$ must lie in the unique torus centralizing
$B$. It exists exactly when $6\mid n$, equivalently $m$ even,
and is then unique. Fixing $M$ and its subgroup $B$ loses no
configuration because both choices form a single chain class.
This proves the first count.

A subgroup $S_3$ containing $B$ normalizes $B$ and is the inverse
image of a reflection subgroup of order two in $K$. Conversely each
reflection $r^k s$, with $k\in\mathbb Z/m\mathbb Z$, gives one
such $S_3$. Indeed its inverse image has order six and inverts $B$.
There are exactly $m$ possibilities. The stabilizer of the chain
$(M,B)$ inside $G$ is $B$, which fixes every one of these overgroups.
Thus all $m$ give distinct conjugacy classes of configurations.
For the $E$-classes we further quotient by $j$, whose action is
$k\mapsto1-k$. This involution has one fixed point for odd $m$
and none for even $m$, proving the count $\lceil m/2\rceil$.

In both types, the intersection with $M$ is exactly $B$:
$B$ is maximal in $M$ and $M$ cannot lie in a group of order six.
Neither group of order six is a subgroup of $A_4$, so $M$ and $H$
generate $G$ by maximality of $M$. Their local actions are
respectively the natural action of $A_4$ of degree four and an
action of degree two. This verifies the rank two geometry and local
two-transitivity conditions, rather than just counting abstract pairs.
\end{proof}

\section{Equal parabolic types and inversion}
\begin{lemma}\label{inversion}
If a group $K$ acts regularly on a set, its orbits on ordered pairs
of distinct points are parametrized by $K\setminus\{1\}$.
Its orbits on unordered pairs are parametrized by the relation
$k\sim k^{-1}$ on that same set.
\end{lemma}
\begin{proof}
Using the regular coordinates, a pair $(a,b)$ has invariant
$a^{-1}b$. Translate its first entry to the identity to see that
this invariant is complete. Reversing the pair replaces the invariant
by its inverse. These are exactly the further identifications when
order is forgotten.
\end{proof}

\begin{theorem}\label{equal}
Under the source's unordered convention, the configurations with two
parabolics $A_4$ meeting in $B$ have
\[
 m+\left\lfloor\frac m2\right\rfloor
 \quad\hbox{conjugacy classes, and}\quad m
 \quad\hbox{classes under }E.
\]
\end{theorem}
\begin{proof}
Two distinct members of $\mathcal A$ meet in exactly $B$, since
$B$ is maximal in each $A_4$. They generate $G$, and both local
actions have degree four and are two-transitive. Conversely every
configuration of the specified type occurs this way. Any conjugacy
between two such pairs normalizes their intersection $B$.
The conjugacy problem is therefore precisely the regular action of
$K$ on unordered pairs in $\mathcal A$.

All $m$ reflections in $K$ are self-inverse. Among the nonidentity
rotations, inversion has $\lfloor m/2\rfloor$ orbits. Lemmas
\ref{regular} and~\ref{inversion} give the asserted inner count.
For $E$, the additional relation is $k\sim\theta(k)$. On rotations
this already is inversion. On reflections inversion is trivial,
whereas $\theta$ is $k\mapsto1-k$ on their exponents. Thus the
number of classes under $E$ is
\[
 \left\lfloor\frac m2\right\rfloor+
 \left\lceil\frac m2\right\rceil=m.
\]
This also shows exactly why the source's full automorphism count
survives the correction to its conjugacy count.
\end{proof}

\section{The two source passages}
Table 2 of the published Part II article has seven rows. Its first
row is Theorem~\ref{different} for $C_6$; combining $m$ even with
the maximal-$A_4$ congruences gives exactly
$p\equiv13,37,83,107\pmod {120}$. Rows 2--5 are the $S_3$
case, separated by the sign of $p-\epsilon$ and the parity of $m$.
All five rows retain their two class counts.

Rows 6--7 concern two $A_4$ parabolics. They print the inner count
$(p-\epsilon)/3-1=2m-1$, which counts ordered pairs in the regular
set. Section 2 of the published source defines configurations using
unordered sets of parabolics. Their correct inner count under that
definition is $m+\lfloor m/2\rfloor$. The printed full automorphism
count $(p-\epsilon)/6=m$ is recovered. For example, the inner count
at $p=37$ is nine instead of eleven, and at $p=43$ it is ten
instead of thirteen. These are consequences of the uniform proof,
not numerical searches.

The fixed appendix repeats the same calculation in Proposition 16,
Subcase 3, on its physical page 4. It subtracts one from the regular
set size but does not then divide by inversion. The same correction
therefore disposes of that count in the second original. Its
Subcases 1--2, physical pages 2--4, have the recovered different-type
counts proved above. The exceptional case $p=5$ has $m=1$ and both
counts equal one, so the correction does not alter that row's value.

The original configurations and the classical subgroup premises are
attributed to the source and its cited subgroup classifications.
The common normalizer description, explicit reflection action and
inversion quotient are the derivation used here. This component does
not complete either original, add a principal refutation, claim a new
open-conjecture answer, claim Lean verification or measure compression.

\begin{thebibliography}{9}
\bibitem{source} F.~Buekenhout, J.~De Saedeleer and D.~Leemans,
\emph{On the rank two geometries of the groups $\PSL(2,q)$: part II},
Ars Mathematica Contemporanea \textbf{6} (2013), 365--388,
Section 2, Table 2 and Proposition 5.6.
\bibitem{appendix} The same authors, \emph{Part II (Appendix)},
retained 21-page edition, Proposition 16, physical pp.~2--4.
\end{thebibliography}
\end{document}
